% ---------------------------------------------------
% Trabajo Final de Grado
% Author: Fabián González Lence <alu0101549491@ull.edu.es>
% Chapter: Conclusions and Future Lines of Work
% ----------------------------------------------------

\selectlanguage{english}
\chapter{Conclusions and Future Lines of Work} \label{chap:Conclusions}

\section{Conclusions} \label{sec:conclusions}

As a result of the work carried out, the \ac{TFG} repository \cite{URL::TFG} contains the code of five deployed and functional web applications whose production code was generated entirely by \ac{AI} models with differentiated roles, capable of producing, in some cases within a single prompt iteration, complete modules of structured, documented code coherent with the design diagrams. Putting this into perspective: in October 2025, when the experiment began with \ac{HANGMAN} and the question was whether Claude could generate a UML diagram in Mermaid, reaching the final result was far from obvious. The results obtained, analysed in Chapter~\ref{chap:Results}, exceed the initial expectations in terms of scale and coherence, although this positive assessment coexists with episodes of mediocre performance, particularly in iterative tasks ---user interface adjustments, quality metrics, among others--- where models tended to require several rounds of instructions to converge on the expected result.

One of the most relevant methodological findings of this work is that the distribution of roles among models proved to be the design decision with the greatest impact on the quality of the final result. As documented in Chapter~\ref{chap:Metodology}, this distribution introduces a layer of abstraction that reproduces the dynamics of a human development team; however, its full benefit did not materialise until the transition ---described in Chapter~\ref{chap:AIModels}--- from conversational \textit{chatbots} external to the \ac{IDE} to custom GitHub Copilot agents integrated directly into the repository. Direct access to the codebase, without the friction of copying and pasting between external tools, enabled a degree of specialisation and decoupling that notably reduced the inter-modular incoherence observed in the projects prior to the change. This shift was not planned; it emerged from the evidence accumulated across the first three projects: the conversational \textit{chatbot} approach is effective for small-scale projects such as \ac{HANGMAN} and \ac{PLAYER}, but does not scale with sufficient coherence when the system grows to the size of \ac{CARTO} or \ac{TENNIS}, with \ac{BALATRO} being the inflection point where its limitations became most evident.

Regarding the human role in the process, the experience confirms that the engineer's differential value lay not in writing code, but in coordinating agents, validating intermediate results, and translating client preferences into \textit{prompt} language ---alternating at all times between the role of technical director and that of client spokesperson---. This function proved indispensable precisely because the models exhibit a recurring limitation: when the project context exceeds their effective window, inconsistent dependencies, duplicated methods, or references to non-existent modules emerge, compromising the overall coherence of the resulting product. Without human intervention as a validation and redirection layer, this inter-modular degradation would have significantly affected the larger-scale projects.

The choice of model at each phase was largely conditioned by this same limitation. Claude ---specifically in its Sonnet 4.5 and 4.6 versions, as analysed in Chapter~\ref{chap:AIModels}--- consistently stood out for combining a large context window with a high capacity to maintain modular coherence and comply with the technical guidelines established in the \textit{prompts}, making it the most valuable model throughout the experiment. Other models, such as Mistral and Qwen, on the other hand, showed notable performance in smaller-scale projects ---especially in generating code and test suites for \ac{HANGMAN} and \ac{PLAYER}---, but degraded significantly in \ac{BALATRO}, where the complexity of the domain model exceeded its capacity to produce functional tests in a stable manner. Beyond these differences between models, the five projects generally present a structure aligned with the class diagrams produced in the design phase and documentation consistent with software engineering standards, demonstrating that the multi-model approach with defined roles, when correctly applied, yields acceptable code quality.

The broadest reflection to emerge from this work concerns the place of \ac{AI} agents in the software development process. Drawing on the experience accumulated across the five projects, it is concluded that, in the current state of the technology, these models do not replace the engineer: they act as an extension of the engineer, analogously to how the emergence of high-level languages did not eliminate the low-level language programmer but transformed the way programmers work. The differential key does not lie in being able to describe to an \ac{AI} what one wishes to build ---that is within anyone's reach---, but in possessing sufficient technical knowledge to evaluate the quality of what is received, identify its deficiencies, and guide the process so that the final result is correct, maintainable, and scalable. The results of this \ac{TFG} ultimately suggest that the \ac{AI}-assisted development paradigm does not eliminate the need for technical training in software engineering, but rather raises the ceiling of what a well-trained engineer can produce in a given time.

\newpage

\section{Future Lines of Work} \label{sec:future-work}

The scale of the experiment, limited to five web projects, constrains the generality of the conclusions obtained. A first avenue for future research would be to increase the number of participants and replicate the experiment with different programmer profiles: asking several engineers to build the same application using the same models and comparing the results would allow isolating and studying the effect of the experimenter's prompting skill from the intrinsic performance of the model.

Complementarily, revisiting the comparative approach discarded in the early stages of this work ---developing the same application in parallel with different models and comparing results in terms of quality, requirements coverage, and supervision effort--- would provide quantitative data on the differences between models that the adopted collaborative approach cannot extract in isolation, and would answer questions that this \ac{TFG} raised but could not address empirically. Extending the experiment to domains beyond web development constitutes another relevant direction: desktop applications, embedded systems, or data-intensive \textit{backends} would introduce different technical constraints that would test the models' capabilities in contexts less represented in their training data and with other programming languages.

Finally, the line of greatest long-term transformative potential is the complete automation of the agent coordination flow: having one \ac{AI} manage other \ac{AI}s without human intervention in directing the process. In the current state of the technology, the role of the human development lead remains the indispensable component of the system; its substitution by an autonomous orchestrator agent would represent a qualitative advance that, although technically plausible, has not yet proven viable in projects of real complexity. Its investigation in the context of \ac{AI}-assisted software development would constitute a high-impact contribution to the field of software engineering.